% Chapter 13: Experiments and Numerical Verification
% Part IV: Implementation and Applications

\chapter{실험과 수치 검증}
\label{ch:experiments}

\begin{quote}
\textit{``수학적 증명은 정리가 참임을 보장하지만, 수치 실험은 정리가 \emph{적절한지}를 보여준다.''}
\end{quote}

\noindent
SCC 이론은 48개의 Category A 정리를 통해 엄밀하게 증명되었다. 그러나 수학적 증명만으로는 답할 수 없는 질문이 있다: 이론적 임계값이 실제로 관찰 가능한 전이를 만드는가? 메타안정성 향상은 얼마나 큰가? 매개변수 민감도는 어떠한가?

이 장에서는 58개의 수치 실험(exp1--exp58)을 통해 이론의 핵심 예측을 검증한다. 모든 실험 코드는 \texttt{experiments/} 디렉토리에서 재현 가능하다. 각 실험은 특정 정리나 예측을 시험하도록 설계되었으며, ``무작위로 돌려본'' 실험은 하나도 없다.


%% ============================================================
\section{실험 설계 철학}
\label{sec:exp-philosophy}

SCC 실험은 다음 원칙을 따른다.

\paragraph{이론-우선 설계.}
모든 실험은 ``어떤 정리의 어떤 주장을 검증하는가?''라는 질문에서 시작한다. 실험 결과가 이론과 불일치하면, 이는 실험의 실패가 아니라 이론의 수정 기회이다. 실제로 12회의 반복 개발(I1--I12) 동안 실험이 이론의 여러 오류를 발견했다:
\begin{itemize}
    \item I6: 이중 우물 기울기의 인수 2 누락 발견
    \item I8: $C_t$-가중 Sep가 형성 품질과 무관하게 $\sim 0.5$를 반환하는 문제 발견
    \item I11: 4-성분 지문에서 레졸벤트 $C(x,x)$의 야코비안 $\sim 9300$ 문제 발견
\end{itemize}

\paragraph{재현성.} 모든 실험은 시드(seed) 42를 사용하는 고정 난수 발생기로 완전히 재현 가능하다. 실행 명령:
\begin{lstlisting}[language=bash]
python3 experiments/exp1_lambda_sweep.py
python3 experiments/exp2_phase_transition.py
# ... 등
\end{lstlisting}


%% ============================================================
\section{매개변수 소거와 민감도 분석}
\label{sec:exp-param-sweeps}

\subsection{$\lambda$-비율 소거 (exp1)}

실험 exp1(\texttt{exp1\_lambda\_sweep.py})은 세 에너지 가중치 $(\lambda_{\mathrm{cl}}, \lambda_{\mathrm{sep}}, \lambda_{\mathrm{bd}})$의 상대적 비율이 형성 품질에 미치는 영향을 조사한다.

\paragraph{설계.} $10 \times 10$ 격자에서 $\lambda_i \in \{0.1, 0.5, 1.0, 2.0, 5.0\}$의 조합을 소거하며, 각 조합에서 진단 벡터를 기록한다.

\paragraph{핵심 발견.}
\begin{itemize}
    \item $\lambda_{\mathrm{bd}} \gg \lambda_{\mathrm{cl}}, \lambda_{\mathrm{sep}}$일 때: Allen-Cahn 유사 형성 (날카로운 경계, 낮은 Bind)
    \item $\lambda_{\mathrm{cl}} \gg$ 나머지: 균일한 장 (폐합 고정점으로 수렴)
    \item $\lambda_{\mathrm{sep}} \gg$ 나머지: 비물리적 분열
    \item 헤시안 정규화 후: 세 항이 균형을 이루며, 진단 벡터의 모든 성분이 $> 0.8$
\end{itemize}

\paragraph{시사점.} 헤시안 스펙트럼 정규화(\S\ref{sec:energy-impl})의 중요성을 확인한다. 정규화 없이는 매개변수 조정이 극히 민감하다.

\subsection{민감도 분석 (exp3, exp5)}

exp3(\texttt{exp3\_ablation.py})은 각 에너지 항을 하나씩 제거하는 절삭(ablation) 실험이고, exp5(\texttt{exp5\_sensitivity.py})는 개별 매개변수의 일차 민감도를 측정한다.

\paragraph{핵심 결과.}
\begin{itemize}
    \item $\mathcal{E}_{\mathrm{cl}}$ 제거: Bind $\downarrow$ 0.3, 나머지 영향 미미 $\to$ 네 에너지 항의 독립성 확인
    \item $\mathcal{E}_{\mathrm{sep}}$ 제거: Sep $\downarrow$ 0.5, Inside 무관 $\to$ 분리와 형태의 독립성
    \item 가장 민감한 매개변수: $\beta_{\mathrm{bd}} / \alpha_{\mathrm{bd}}$ 비율 (위상 전이 임계값 결정)
\end{itemize}

% Figure placeholder
\begin{figure}[t]
\centering
\fbox{\parbox{0.85\textwidth}{\centering\vspace{5cm}
\textbf{[그림 13.1]} $\lambda$-비율 소거 결과.\\
좌: 정규화 전, 진단 벡터 성분의 넓은 분산.\\
우: 헤시안 정규화 후, 네 성분이 $0.8$--$1.0$ 범위에 안정화.
\vspace{0.5cm}}}
\caption{에너지 가중치 소거 (exp1). 헤시안 스펙트럼 정규화가 매개변수 공간에서의 강건성을 크게 향상시킨다.}
\label{fig:lambda-sweep}
\end{figure}


%% ============================================================
\section{위상 전이 검증}
\label{sec:exp-phase-transition}

T8-Core의 핵심 예측은 $\beta / \alpha > 4\lambda_2 / |W''(c)|$일 때 비균일 최소화자가 존재한다는 것이다. 세 실험이 이를 다각도로 검증한다.

\subsection{위상 전이 위치 (exp2)}

exp2(\texttt{exp2\_phase\_transition.py})는 $\beta$를 연속적으로 변화시키며 에너지 차 $\Delta\mathcal{E} = \mathcal{E}(\hat{u}) - \mathcal{E}(c \cdot \mathbf{1})$를 추적한다.

\paragraph{결과.} $\beta = \beta_{\mathrm{crit}}$ 부근에서 급격한 전이: $\Delta\mathcal{E}$가 0에서 음수로 전환된다. 이론적 $\beta_{\mathrm{crit}}$와의 일치도 $R^2 > 0.99$.

\subsection{분기점 교차와 히스테리시스 부재 (exp37)}

T-Birth-Parametric은 초임계 갈림길 분기(supercritical pitchfork)를 예측한다. exp37(\texttt{exp37\_bifurcation\_crossing.py})은 $\beta$를 $\beta_{\mathrm{crit}}$ 전후로 순방향/역방향 소거하여 히스테리시스의 부재를 확인한다.

\paragraph{핵심 발견.}
\begin{itemize}
    \item 순방향 ($\beta$ 증가) 전이점과 역방향 ($\beta$ 감소) 전이점이 일치: \textbf{히스테리시스 없음}
    \item 비균일 최소화자의 진폭이 $|\hat{u} - c| \propto (\beta - \beta_{\mathrm{crit}})^{1/2}$로 성장: 초임계 분기 확인
\end{itemize}

\subsection{위상적 형성 탄생 (exp39)}

exp39(\texttt{exp39\_formation\_birth.py})는 $\beta_{\mathrm{crit}}$ 이전(평탄 에너지 지형)과 이후(형성 출현)의 에너지 지형을 비교한다.

\paragraph{결과.} $\beta < \beta_{\mathrm{crit}}$에서는 모든 다중 시작이 균일 상태로 수렴. $\beta > \beta_{\mathrm{crit}}$에서는 Fiedler 고유벡터 방향으로 비균일 형성이 자발적으로 출현한다.

\subsection{스펙트럼 보편성 (FORMATION-BIRTH, 32개 그래프)}

FORMATION-BIRTH 정리는 $\beta_{\mathrm{crit}}$가 $\lambda_2$ 하나로만 결정된다는 스펙트럼 보편성을 주장한다. exp2와 추가 실험에서 32개의 다양한 그래프(격자, 바벨, SBM, 무작위 기하 그래프)에 대해 이론적 $\beta_{\mathrm{crit}}$와 관측 전이점의 상관을 측정했다.

\paragraph{결과.} $R^2 = 0.9924$. $\lambda_2$가 그래프 위상에 무관하게 위상 전이를 지배한다.

% Figure placeholder
\begin{figure}[t]
\centering
\fbox{\parbox{0.85\textwidth}{\centering\vspace{5cm}
\textbf{[그림 13.2]} 위상 전이 검증.\\
좌: $\beta$ 대 에너지 차 $\Delta\mathcal{E}$ (exp2). 이론적 $\beta_{\mathrm{crit}}$ 위치(점선)에서 전이 발생.\\
우: 32개 그래프의 스펙트럼 보편성 산점도 ($R^2 = 0.9924$).
\vspace{0.5cm}}}
\caption{위상 전이 검증. 이론적 임계값 $\beta_{\mathrm{crit}} = 4\alpha\lambda_2/|W''(c)|$이 다양한 그래프 위상에서 정확히 성립한다.}
\label{fig:phase-transition}
\end{figure}


%% ============================================================
\section{Allen-Cahn 비교}
\label{sec:exp-allen-cahn}

\subsection{메타안정성 향상 (exp7)}

exp7(\texttt{exp7\_allen\_cahn\_vs\_scc.py})은 순수 Allen-Cahn 에너지($\mathcal{E}_{\mathrm{bd}}$만)와 SCC 전체 에너지를 비교한다.

\paragraph{실험 설계.}
\begin{enumerate}
    \item 동일 격자와 $\alpha, \beta$에서 양쪽 모두 형성을 찾음
    \item 수렴된 형성의 헤시안 최소 고유값 $\mu_{\min}$을 비교
    \item 노이즈($\sigma = 0.01$--$0.10$)에 대한 형성 존속 시간 측정
\end{enumerate}

\paragraph{핵심 결과.}
\begin{itemize}
    \item SCC 형성의 $\mu_{\min}$이 Allen-Cahn보다 일관되게 높음 (T7-Enhanced 확인)
    \item 에너지 장벽 차이: SCC 형성이 노이즈 하에서 $1.5$--$3\times$ 더 오래 존속
    \item SCC의 추가 이점: Sep $> 0.9$ (Allen-Cahn은 Sep 개념 없음)
\end{itemize}

\paragraph{물리적 해석.} 폐합 항 $\mathcal{E}_{\mathrm{cl}}$이 폐합 고정점 주변에 양정치 헤시안 기여를 추가하여 (Gram 행렬 $2(I - J_{\mathrm{Cl}})^T(I - J_{\mathrm{Cl}})$), 비단순 위치에서의 안정성을 강화한다. 이것이 T3/T6-Stability의 계산적 실현이다.

% Figure placeholder
\begin{figure}[t]
\centering
\fbox{\parbox{0.85\textwidth}{\centering\vspace{5cm}
\textbf{[그림 13.3]} Allen-Cahn 대 SCC 비교 (exp7).\\
좌: 헤시안 최소 고유값 비교---SCC가 일관되게 더 높음.\\
우: 노이즈 하 형성 존속 시간---SCC 형성이 1.5--3$\times$ 더 오래 지속.
\vspace{0.5cm}}}
\caption{SCC의 메타안정성 향상 (exp7). 자기-참조적 폐합 보정이 에너지 장벽을 깊게 한다.}
\label{fig:allen-cahn-comparison}
\end{figure}


%% ============================================================
\section{전달과 지속성}
\label{sec:exp-transport}

\subsection{전달 집중 (exp10--11)}

exp10(\texttt{exp10\_fingerprint\_verification.py})은 응집 지문 $\varphi = (u, \mathrm{Cl}(u), D(x;1{-}u))$의 판별력을 검증하고, exp11(\texttt{exp11\_transport\_concentration.py})은 핵심-대-핵심 전달 집중을 측정한다.

\paragraph{핵심 결과.}
\begin{itemize}
    \item 심층 핵심($\delta \geq 2$) 지문 갭: $\Delta_\varphi^2 = 2.44$ (이론: 2.87, 같은 차수)
    \item 핵심-대-핵심 전달 분율: $> 99.99\%$ ($\gamma / \varepsilon_{\mathrm{OT}} > 5$에서)
    \item 약한 레짐 고정점: 2--3회 반복으로 수렴
\end{itemize}

\subsection{전달 체인 검증 (exp26--27)}

T-Persist-1의 전체 5개 구성요소 $(a)$--$(e)$를 검증하는 종단간(end-to-end) 실험이다.

exp26(\texttt{exp26\_full\_chain\_verification.py})은 표준 다중 시작에서 $(a)(c)(e)$가 보편적으로 통과하지만, $(b)(d)$가 분지 전환(basin-switching)으로 실패할 수 있음을 보인다. exp27(\texttt{exp27\_warm\_start\_chain.py})은 웜-스타트(warm-start)를 사용하여 \textbf{5/5 구성요소 $\times$ 5/5 구성 = 100\% 통과}를 달성한다.

\subsection{스트레스 테스트 (exp28)}

exp28(\texttt{exp28\_stress\_test.py})은 100개의 매개변수-격자 조합에서 전달 체인의 강건성을 테스트한다.

\paragraph{결과.} $\varepsilon \leq 0.20$에서 84/100 통과. 16개 실패는 모두 소규모 격자($n < 64$) 또는 약한 위상 분리($\beta < 20$)에서 발생. $n \geq 64, \beta \geq 20$에서는 모든 조합이 통과한다.

% Figure placeholder
\begin{figure}[t]
\centering
\fbox{\parbox{0.85\textwidth}{\centering\vspace{5cm}
\textbf{[그림 13.4]} 전달과 지속성 검증.\\
좌상: 지문 갭 $\Delta_\varphi^2$ 분포 (exp10).\\
우상: 핵심-대-핵심 전달 분율 대 $\gamma/\varepsilon_{\mathrm{OT}}$ (exp11).\\
하: 스트레스 테스트 84/100 통과 매트릭스 (exp28)---실패는 소격자/약분리에 집중.
\vspace{0.5cm}}}
\caption{전달 집중 및 지속성 체인 검증. 이론의 조건 (H2', H3, GT)이 만족되면 5/5 구성요소가 보편적으로 통과한다.}
\label{fig:transport-verification}
\end{figure}


%% ============================================================
\section{다중 형성 실험}
\label{sec:exp-multi}

\subsection{$K$-형성 최적화 (exp6)}

exp6(\texttt{exp6\_multiformation.py})은 $K = 2, 3, 4$ 형성을 다양한 격자에서 찾는다. 형성 간 반발($\lambda_{\mathrm{rep}}$)이 공간적 분리를 유도하고, 심플렉스 장벽이 총 참여를 1 이하로 제한함을 확인한다.

\subsection{합병 흐름과 에너지 장벽 (exp30, exp38)}

T-Merge 정리에 따르면 $K > 1$ 형성은 메타안정 지역 최소점이며, $K = 1$이 전역 최소이다.

exp30(\texttt{exp30\_merge\_flow.py})은 두 형성을 점진적으로 접근시키며 합병 전이를 관찰한다. exp38(\texttt{exp38\_barrier\_height.py})은 합병 에너지 장벽 $\Delta E_{\mathrm{merge}}$를 측정한다.

\paragraph{핵심 결과.}
\begin{itemize}
    \item 장벽 스케일링: $\Delta E \propto \beta^{1.24}$ (이론: $\beta^{0.89}$). 실제 장벽이 이론 하한보다 높음---보수적 한계 확인
    \item 합병 전이는 형성 간 거리 $d_{\min}$이 임계거리 이하로 떨어질 때 발생
\end{itemize}

\subsection{확률적 조대화 (exp55)}

exp55(\texttt{exp55\_stochastic\_coarsening.py})는 노이즈 하에서 $K \to K-1$ 전이의 Kramers 장벽 교차 동역학을 시뮬레이션한다. $K > 1$ 형성의 수명은 $\propto e^{\Delta E / T}$로 노이즈 온도 $T$에 지수적으로 의존한다.

\subsection{레짐 분류 검증 (exp46--49)}

T-Persist-K-Unified의 $\Lambda_{\mathrm{coupling}}$ 기반 레짐 분류를 69개 구성에서 검증한다.

\paragraph{결과.} 기하학적 $\Lambda$와 스펙트럼 분석의 일치도 100\%. 세 레짐(잘-분리, 약한 상호작용, 강한 상호작용)이 예측대로 구분된다.

% Figure placeholder
\begin{figure}[t]
\centering
\fbox{\parbox{0.85\textwidth}{\centering\vspace{5cm}
\textbf{[그림 13.5]} 다중 형성 실험.\\
좌상: $K=2$ 형성의 에너지 지형 (exp6).\\
우상: 합병 에너지 장벽 대 $\beta$ (exp38), $\Delta E \propto \beta^{1.24}$.\\
하: $\Lambda_{\mathrm{coupling}}$ 위상 도표---세 레짐의 경계가 이론적 임계값과 일치 (exp47).
\vspace{0.5cm}}}
\caption{다중 형성 실험 결과. $K > 1$ 형성의 메타안정성과 레짐 분류가 이론적 예측과 일치한다.}
\label{fig:multi-formation-exps}
\end{figure}


%% ============================================================
\section{스케일링과 대규모 격자}
\label{sec:exp-scaling}

\subsection{스케일 검증 (exp42)}

exp42(\texttt{exp42\_scale\_verification.py})는 $10 \times 10$에서 $30 \times 30$까지 격자 크기를 변화시키며, 형성 품질(진단 벡터)과 계산 비용을 측정한다.

\paragraph{결과.}
\begin{itemize}
    \item 진단 벡터 품질은 격자 크기에 대해 안정적 (Bind, Sep $> 0.9$)
    \item Inside는 격자가 커질수록 약간 향상 (깊은 핵심 비율 증가)
    \item 계산 비용: 희소 행렬 연산으로 $O(n)$ 수준의 반복당 비용
\end{itemize}

\subsection{$50 \times 50$ 스케일 테스트 (exp43)}

exp43(\texttt{exp43\_50x50\_scale.py})은 2500 노드 격자에서 SCC 파이프라인의 작동을 확인한다.

\paragraph{결과.} $50 \times 50$에서도 형성 발견 성공. 2500 노드에서 수렴까지 약 200--500회 반복. 반-암묵적 방식이 순수 명시적 방식 대비 $3\times$--$5\times$ 적은 반복으로 수렴한다.

% Figure placeholder
\begin{figure}[t]
\centering
\fbox{\parbox{0.85\textwidth}{\centering\vspace{5cm}
\textbf{[그림 13.6]} 스케일링 결과 (exp42--43).\\
좌: 격자 크기 대 진단 벡터 성분---품질 안정적.\\
우: 격자 크기 대 수렴 반복 수---반-암묵적 방식의 효율성.
\vspace{0.5cm}}}
\caption{스케일링 분석. SCC 파이프라인은 2500 노드까지 안정적으로 작동하며, 반-암묵적 방식이 수렴 속도를 유의하게 향상시킨다.}
\label{fig:scaling}
\end{figure}


%% ============================================================
\section{예측 검증}
\label{sec:exp-predictions}

SCC 이론은 다섯 가지 계산적 예측(P1--P5)을 제시한다. 이들은 다른 모든 정리와 독립적으로 검증 가능하다.

\subsection{다섯 예측과 검증 결과 (exp15)}

exp15(\texttt{exp15\_prediction\_verify.py})는 다섯 예측을 일괄 검증한다.

\begin{table}[h]
\centering
\caption{다섯 계산적 예측의 검증 결과}
\label{tab:predictions}
\begin{tabular}{clcl}
\toprule
\textbf{예측} & \textbf{내용} & \textbf{결과} & \textbf{비고} \\
\midrule
P1 & 폐합에 의한 축약 & \checkmark & $\mathrm{Cl}^k(u) \to u^*$ 관측 \\
P2 & 에너지 항 독립성 & \checkmark & 절삭 실험으로 확인 (exp3) \\
P3 & 향상된 체류 시간 & \checkmark & $p = 0.037$ (Mann-Whitney) \\
P4 & 경로 의존성 & \checkmark & 다중 시작에서 다른 최소점 \\
P5 & Sep-before-Inside 순서 & \checkmark & Sep 수렴이 Inside보다 선행 \\
\bottomrule
\end{tabular}
\end{table}

\paragraph{P3: 향상된 체류 시간.} 가장 비자명한 예측이다. SCC 형성이 Allen-Cahn 형성보다 노이즈 하에서 더 오래 존속한다는 주장으로, Mann-Whitney 검정에서 $p = 0.037$으로 통계적으로 유의미하다.

\paragraph{P5: Sep-before-Inside 순서.} 최적화 과정에서 Sep 진단이 Inside 진단보다 먼저 높은 값에 도달하는 현상이다. 이는 구별(내부/외부 비대칭)이 형태적 분절보다 에너지적으로 먼저 달성되기 때문이다.

\subsection{Bind-$\tau$ 소거와 일반 $\tau$ 검증 (exp58)}

exp58(\texttt{exp58\_bind\_tau\_sweep.py})은 T-Bind-Proj/Full의 일반 $\tau_{\mathrm{cl}}$ 확장을 검증한다. 68개의 $(\tau, a_{\mathrm{cl}}, c)$ 조합에서 이진 질량-균형 공식 $\Phi(\tau; a_{\mathrm{cl}}, c)$과의 일치도 $R^2 = 0.995$를 확인한다.

\paragraph{핵심 발견.}
\begin{itemize}
    \item $\tau = 1/2$ (기본값)에서 $\bar{r}_0 = O(n^{-1/d})$: 이진 근사가 정확
    \item $\tau \neq 1/2$에서 $\bar{r}_0 = O(1)$: 폐합 고정점의 비대칭성으로 질량-균형 상쇄가 실패
    \item 체적-호환 문턱값 $\tau^*(c)$의 존재 확인
\end{itemize}

이 결과는 정리 6.1(원래 $c \neq 1/2$에 대한 KKT 논증)의 철회와 보정된 정리 6.1'의 정당성을 실험적으로 뒷받침한다.

% Figure placeholder
\begin{figure}[t]
\centering
\fbox{\parbox{0.85\textwidth}{\centering\vspace{5cm}
\textbf{[그림 13.7]} 예측 검증 결과.\\
좌: P3 체류 시간 분포---SCC(파랑) vs. Allen-Cahn(주황), $p = 0.037$.\\
우: P5 Sep-before-Inside 순서---최적화 반복에 따른 진단 성분 궤적.
\vspace{0.5cm}}}
\caption{다섯 계산적 예측의 검증. P3의 향상된 체류 시간과 P5의 Sep 선행 순서가 특히 비자명하다.}
\label{fig:predictions}
\end{figure}

% Figure placeholder for Bind-tau
\begin{figure}[t]
\centering
\fbox{\parbox{0.85\textwidth}{\centering\vspace{5cm}
\textbf{[그림 13.8]} Bind-$\tau$ 소거 (exp58).\\
$x$축: 이론적 $\bar{r}_0(\tau)$, $y$축: 관측 $\bar{r}_0$.\\
$R^2 = 0.995$, 68개 구성. $\tau = 1/2$ 부근에서 $\bar{r}_0 \to 0$ (격자 크기 증가 시).
\vspace{0.5cm}}}
\caption{일반 $\tau_{\mathrm{cl}}$에 대한 T-Bind-Proj 검증. 이진 질량-균형 공식이 넓은 매개변수 범위에서 정확하다.}
\label{fig:bind-tau}
\end{figure}


%% ============================================================
\section{고급 실험: 레짐 소거에서 핵형성까지}
\label{sec:exp-advanced}

이 절에서는 이론 개발의 후반부(Phase 9--14)에서 수행된 고급 실험을 요약한다. 이 실험들은 T-Persist-K-Unified, 스펙트럼 $K$-선택, 폐합 문턱값 분석 등 이론의 가장 정교한 부분을 검증한다.

\subsection{레짐 소거와 통합 예측 (exp46--49)}

T-Persist-K-Unified 정리는 세 레짐(잘-분리, 약한 상호작용, 강한 상호작용)을 단일 매개변수 $\Lambda_{\mathrm{coupling}}$으로 통합한다. exp46--49는 이 통합의 정량적 정확성을 검증한다.

\begin{itemize}
    \item \texttt{exp46\_weak\_to\_strong\_transition.py}: $\Lambda$ 증가에 따른 $\mathsf{Persist}$ 열화---이론적 $1 - C\Lambda^2$ 예측과의 일치
    \item \texttt{exp47\_phase\_diagram.py}: $(\lambda_{\mathrm{rep}}, d_{\min})$ 평면의 레짐 경계---69개 구성에서 100\% 일치
    \item \texttt{exp48\_regime\_sweep.py}: $K = 2, 3, 4$에서의 레짐 분류 일관성
    \item \texttt{exp49\_unified\_predictions.py}: 통합 정리의 네 가지 결론 동시 검증
\end{itemize}

\paragraph{핵심 발견.} $\Lambda_{\mathrm{coupling}} < 1/(K-1)$ 조건이 지속성을 보장하는 것이 모든 구성에서 확인되었다. $\Lambda > 1/(K-1)$에서는 합병 분기가 발생하며, 이는 T-Merge의 예측과 일치한다.

\subsection{스펙트럼 $K$-선택 (exp51)}

exp51(\texttt{exp51\_k\_selection.py})은 라플라시안 스펙트럼으로부터 최적 형성 수 $K$를 예측한다. 위상 전이 문턱값 이하의 고유값 수가 $K^*$를 예측한다:
\begin{equation}
    K^* = |\{k : \lambda_k < \beta_{\mathrm{crit}}\}|
\end{equation}

\paragraph{결과.} $8 \times 8$--$20 \times 20$ 격자에서 예측 $K^*$와 실제 최적 $K$의 일치율 85\%. 불일치는 $\lambda_k \approx \beta_{\mathrm{crit}}$ 경계 근처에서 발생한다.

\subsection{핵형성과 폐합 문턱값 (exp54, exp56, exp57)}

\begin{itemize}
    \item \texttt{exp54\_closure\_threshold.py}: $\tau_{\mathrm{cl}}$의 체적-호환 문턱값 $\tau^*(c)$에서 Bind가 최대화됨을 확인
    \item \texttt{exp56\_nucleation.py}: $K > 1$ 형성의 핵형성 과정---Fiedler 방향 초기화가 필수
    \item \texttt{exp57\_closure\_thorough.py}: $a_{\mathrm{cl}} \in [1, 3.9]$ 전체 범위에서 폐합 축약과 형성 품질 관계
\end{itemize}

\paragraph{핵심 발견.} $a_{\mathrm{cl}} = 3.5$ (기본값)이 Bind와 수렴 속도의 최적 균형점이다. $a_{\mathrm{cl}} \to 4$에서 축약이 약해지며 수렴이 느려진다. $a_{\mathrm{cl}} < 2$에서는 폐합이 너무 약하여 Bind $< 0.7$.

\subsection{전체 실험 목록}

완결성을 위해 58개 실험의 전체 목록을 제공한다:

\begin{table}[h]
\centering
\caption{58개 실험 완전 목록}
\label{tab:all-experiments}
\small
\begin{tabular}{rll}
\toprule
\textbf{\#} & \textbf{파일명} & \textbf{검증 대상} \\
\midrule
1 & \texttt{exp1\_lambda\_sweep} & $\lambda$-비율 민감도 \\
2 & \texttt{exp2\_phase\_transition} & T8-Core $\beta_{\mathrm{crit}}$ \\
3 & \texttt{exp3\_ablation} & 에너지 항 절삭 \\
5 & \texttt{exp5\_sensitivity} & 매개변수 일차 민감도 \\
6 & \texttt{exp6\_multiformation} & $K$-형성 기본 \\
7 & \texttt{exp7\_allen\_cahn\_vs\_scc} & T7-Enhanced \\
8 & \texttt{exp8\_computational\_predictions} & P1--P5 예비 \\
9 & \texttt{exp9\_temporal\_transport} & 전달 커널 기본 \\
10 & \texttt{exp10\_fingerprint\_verification} & 지문 판별력 \\
11 & \texttt{exp11\_transport\_concentration} & 핵심-대-핵심 집중 \\
12 & \texttt{exp12\_multi\_temporal} & $K$-형성 시간 전달 \\
13 & \texttt{exp13\_core\_depth\_verification} & 심층 핵심 존재 \\
14 & \texttt{exp14\_multi\_persist} & 다중 지속성 \\
15 & \texttt{exp15\_prediction\_verify} & P1--P5 최종 \\
16--17 & \texttt{exp16--17} & 전이층, 에너지 장벽 \\
18--25 & \texttt{exp18--25} & 핵심 깊이, 안장, 분지 \\
26--28 & \texttt{exp26--28} & 전달 체인/스트레스 \\
29--36 & \texttt{exp29--36} & 전달 소거, 합병, 방향 \\
37--41 & \texttt{exp37--41} & 분기 교차, 가두리 \\
42--43 & \texttt{exp42--43} & 스케일링 검증 \\
44--49 & \texttt{exp44--49} & 종합 검증, 레짐 \\
50--53 & \texttt{exp50--53} & KKT, $K$-선택, 이완 \\
54--58 & \texttt{exp54--58} & 폐합, 핵형성, Bind-$\tau$ \\
\bottomrule
\end{tabular}
\end{table}

% Figure placeholder
\begin{figure}[t]
\centering
\fbox{\parbox{0.85\textwidth}{\centering\vspace{5cm}
\textbf{[그림 13.9]} 레짐 위상 도표 (exp47).\\
$x$축: $d_{\min}$ (형성 간 거리), $y$축: $\lambda_{\mathrm{rep}}$ (반발 강도).\\
세 레짐의 경계가 $\Lambda_{\mathrm{coupling}}$ 등고선으로 정확히 구분된다.
\vspace{0.5cm}}}
\caption{$\Lambda_{\mathrm{coupling}}$ 기반 레짐 위상 도표. 69개 구성에서 이론적 레짐 분류와 100\% 일치한다.}
\label{fig:regime-phase-diagram}
\end{figure}


%% ============================================================
\section{연습 문제}
\label{sec:ch13-exercises}

\begin{enumerate}
    \item \textbf{자신만의 실험 설계.} ``폐합 반복 횟수가 형성 품질에 미치는 영향''을 검증하는 실험을 설계하라. 어떤 정리/예측을 검증하는가? 제어 변수와 독립/종속 변수는 무엇인가?

    \item \textbf{결과 재현.} exp2의 코드를 읽고 $15 \times 15$ 격자에서 실행하라. 이론적 $\beta_{\mathrm{crit}}$를 계산하고, 관찰된 전이점과 비교하라.

    \item \textbf{Allen-Cahn 비교 확장.} exp7을 변형하여 $\mathcal{E}_{\mathrm{sep}}$만 추가한 경우 (폐합 없음)와 $\mathcal{E}_{\mathrm{cl}}$만 추가한 경우 (분리 없음)의 메타안정성을 비교하라. 어떤 항이 메타안정성에 더 큰 기여를 하는가?

    \item \textbf{스트레스 테스트 확장.} exp28의 100개 구성에 $a_{\mathrm{cl}} = 3.0, 3.9$를 추가하여 200개 구성으로 확장하라. 축약 강도가 지속성에 미치는 영향을 분석하라.

    \item \textbf{비격자 그래프.} Erd\H{o}s-R\'{e}nyi 무작위 그래프 $G(50, 0.1)$에서 형성을 찾아라. 격자와 비교하여 진단 벡터가 어떻게 달라지는가? $\lambda_2$가 어떤 역할을 하는가?

    \item \textbf{예측 P3 재검증.} exp15의 체류 시간 실험을 반복하되, 노이즈 수준을 $\sigma = 0.01, 0.05, 0.10, 0.20$으로 변화시켜라. Mann-Whitney $p$-값이 어떻게 변하는가?
\end{enumerate}


%% ============================================================
\section{요약: 58개 실험의 지형도}
\label{sec:exp-summary}

\begin{table}[h]
\centering
\caption{실험 분류표 (주요 실험 요약)}
\label{tab:experiment-summary}
\small
\begin{tabular}{llll}
\toprule
\textbf{범주} & \textbf{실험} & \textbf{검증 대상} & \textbf{결과} \\
\midrule
매개변수 & exp1, 3, 5 & $\lambda$ 민감도, 절삭 & 정규화 효과 확인 \\
위상 전이 & exp2, 37, 39 & T8-Core, T-Birth & $R^2 > 0.99$ \\
Allen-Cahn & exp7 & T7-Enhanced & $1.5$--$3\times$ 향상 \\
전달 & exp9--11, 26--28 & T-Persist-1 & 100\% (웜-스타트) \\
다중 형성 & exp6, 30, 38, 55 & T-Merge, 장벽 & $\beta^{1.24}$ 스케일링 \\
스케일링 & exp42--43 & 대규모 격자 & $50 \times 50$ 성공 \\
레짐 & exp46--49 & T-Persist-K-Unified & 100\% 일치 \\
예측 & exp15, 58 & P1--P5, Bind-$\tau$ & 5/5 통과 \\
\bottomrule
\end{tabular}
\end{table}

58개 실험의 공통된 결론은 명확하다: SCC 이론의 수학적 예측은 계산적으로 정확하며, 이론의 한계(소격자 심층 핵심 부재, 약한 위상 분리에서의 지속성 실패)도 이론 자체의 조건으로 정확하게 예측된다.

다음 장에서는 이 이론이 인지과학, 컴퓨터 비전, 신경과학과 어떻게 연결되는지를 탐구한다.
